\cleardoublepage
\phantomsection

\vspace*{-2.2cm}

{\noindent\Huge\bfseries Resumen\par}
\vspace{0.05cm}
\noindent\rule{\textwidth}{1pt}

\vspace{0.4cm}

El riego en pequeños cultivos, huertos experimentales y espacios agrícolas de baja escala suele depender todavía de rutinas manuales, temporizadores fijos o supervisión presencial. Este enfoque puede resultar poco flexible cuando las condiciones ambientales cambian, cuando no existe presencia continua o cuando se desea conocer con mayor precisión qué está ocurriendo en el suelo y en el entorno de cultivo. En este contexto, TierraViva propone un prototipo de riego inteligente, estable y sostenible basado en Internet de las Cosas (IoT), orientado a combinar sensorización ambiental, decisión local de riego, comunicación móvil y supervisión web.

El sistema desarrollado se organiza como una arquitectura distribuida formada por dos estaciones de campo, una plataforma intermedia de telemetría, una Raspberry Pi 5 como nodo local de supervisión y un \textit{dashboard} web. Cada estación integra sensores de humedad del suelo, temperatura, humedad ambiental, presión y luminosidad, junto con un módulo LTE A7670E-LASA, un relé y una bomba de agua. La decisión de riego no se delega en la nube ni en la interfaz web, sino que se ejecuta localmente en el \textit{firmware} de la estación mediante reglas verificables, modo seguro, confirmación de sequedad, actuación temporizada y periodo de \textit{cooldown}. Esta decisión arquitectónica permite separar claramente el control físico del sistema de las tareas de supervisión y refuerza la estabilidad operativa del prototipo ante fallos de conectividad o disponibilidad de servicios externos.

Las estaciones publican telemetría y códigos de estado en ThingSpeak mediante LTE/HTTP. A partir de esos datos, la Raspberry Pi consulta la información disponible, la normaliza, la almacena en SQLite y la expone mediante una API desarrollada con FastAPI. El \textit{dashboard} permite visualizar lecturas actuales, históricos, gráficas, eventos de sistema y estado de cada estación, proporcionando una capa de supervisión comprensible y accesible.

El prototipo se ha validado mediante pruebas por bloques, registros de monitor serie, publicación en ThingSpeak, integración con Raspberry Pi, visualización web y una prueba funcional en Agrolab Móstoles con la estación B en un entorno real de cultivo. Aunque el trabajo no pretende demostrar todavía un ahorro hídrico cuantificado ni una validación agronómica prolongada, sí alcanza una integración funcional de extremo a extremo y demuestra una base técnica modular, ampliable y defendible para futuras mejoras en riego inteligente y agricultura sostenible.

\textbf{Palabras clave:} Internet de las Cosas, riego inteligente, agricultura sostenible, agricultura de precisión, sensorización ambiental, LTE, ThingSpeak, Raspberry Pi, FastAPI, SQLite, \textit{dashboard}, decisión local de riego.